\documentclass[12pt,letterpaper]{report}
\usepackage{graphicx}
\usepackage{scrextend}
\usepackage{vmargin}
\usepackage{graphicx}
\usepackage{multirow}
\usepackage[utf8]{inputenc}
\usepackage[spanish]{babel}
\usepackage{multicol}
\usepackage{enumerate}
\usepackage{float}
\usepackage{lipsum}
\usepackage{ragged2e}
\usepackage{amsmath, amsthm, amssymb, amsfonts}
\usepackage[usenames]{color}
\parindent=0mm
\pagestyle{empty}
\definecolor{miorange}{rgb}{0.91, 0.43, 0.0}
\begin{document}
\setmargins{2.5cm}      
{1.5cm}                     
{2cm}  
{24cm}                    
{10pt}                          
{1cm}                          
{0pt}                             
{2cm}
\begin{flushright}
\today
\end{flushright}
\vspace{-1.75cm}
\section*{Tarea Clase 7}
\subsection*{Obtenga las soluciones de las siguientes ecuaciones:}
\begin{enumerate}
    \item\begin{equation*}
        \left\lbrace \begin{matrix}
             4x+3y=10 \\
             5x-6y=-7
        \end{matrix} \right.
    \end{equation*}
    \item\begin{equation*}
        \left\lbrace \begin{matrix}
           2x+y=-10 \\
           x-3y=2  
        \end{matrix} \right.
    \end{equation*}
    \item\begin{equation*}
        \left\lbrace \begin{matrix}
             6r-5t=-11\\
             7t-8r=15
        \end{matrix} \right.
    \end{equation*}
    \item\begin{equation*}
        \left\lbrace \begin{matrix}
             12u-16v=24\\
             3u-4v=6
        \end{matrix} \right.
    \end{equation*}
    \item\begin{equation*}
        \left\lbrace \begin{matrix}
             2x+y=9\\
             8x+4y=36
        \end{matrix} \right.
    \end{equation*}
    \item\begin{equation*}
        \left\lbrace \begin{matrix}
             4p-3q=-2\\
             20p-15q=-1
        \end{matrix} \right.
    \end{equation*}
    \item\begin{equation*}
        \left\lbrace \begin{matrix}
             6x-7y=5\\
             8x-9y=7
        \end{matrix} \right.
    \end{equation*}
    \item\begin{equation*}
        \left\lbrace \begin{matrix}
             x+5y=5\\
             3x-5y=3
        \end{matrix} \right.
    \end{equation*}
    \item\begin{equation*}
        \left\lbrace \begin{matrix}
             y=-2x+1\\
             4x+2y=3
        \end{matrix} \right.
    \end{equation*}
    \item\begin{equation*}
        \left\lbrace \begin{matrix}
             2y-3x=1\\
             -4y+6x=-2
        \end{matrix} \right.
    \end{equation*}
    \item\begin{equation*}
        \left\lbrace \begin{matrix}
             6x-5y=-3\\
             3x+2y=12
        \end{matrix} \right.
    \end{equation*}
\end{enumerate}
Las respuestas de las tareas pueden ser enviadas al siguiente correo giovannilopez9808@gmail.com con el asunto entre corchetes de la siguiente manera [Tarea \#] y el nombre del archivo con sus apellidos y nombre, ya sea el nombre completo o solo una parte, por ejemplo \textit{LopezPadilla.pdf}, \textit{LopezGiovanni.pdf}, \textit{LopezPadillaGiovanniGamaliel.pdf}. Las respuestas envienlas en formato pdf, estas pueden ser totalmente digitales (que escriban todo en la computadora) o que sean fotografías, pero por favor, envielo en solo un documento.
\end{document}